% ────────────────────────────────────────────────────────
\section{Derivadas Parciais e Diferenciabilidade}

% ────────────────────────────────────────────────────────
\subsection{Enquadramento Teórico}

\begin{defbox}[title={Derivada Parcial}]
  A \emph{derivada parcial} de $f$ em ordem a $x$ no ponto $(a,b)$ é
  \[
    \frac{\partial f}{\partial x}(a,b)
    = \lim_{h\to 0}\frac{f(a+h,b)-f(a,b)}{h},
  \]
  calculada derivando $f$ em ordem a $x$ mantendo $y$ constante.
\end{defbox}

\begin{defbox}[title={Gradiente}]
  O \emph{gradiente} de $f$ em $(a,b)$ é o vetor
  \[
    \nabla f(a,b) = \left(\frac{\partial f}{\partial x}(a,b),\,
                          \frac{\partial f}{\partial y}(a,b)\right).
  \]
  Aponta na direção de maior crescimento de $f$.
\end{defbox}

\begin{defbox}[title={Derivada Direcional}]
  Para um vetor unitário $\mathbf{u}=(u_1,u_2)$, a derivada direcional é
  \[
    D_{\mathbf{u}}f(a,b) = \nabla f(a,b)\cdot\mathbf{u}
      = \frac{\partial f}{\partial x}u_1 + \frac{\partial f}{\partial y}u_2.
  \]
\end{defbox}

\begin{defbox}[title={Diferenciabilidade}]
  $f$ é \emph{diferenciável} em $(a,b)$ se existir uma aplicação linear
  $L:\mathbb{R}^2\to\mathbb{R}$ tal que
  \[
    \lim_{(h,k)\to(0,0)}\frac{f(a+h,b+k)-f(a,b)-L(h,k)}{\sqrt{h^2+k^2}}=0.
  \]
  A aplicação linear é $L(h,k)=\nabla f(a,b)\cdot(h,k)$.
\end{defbox}

\begin{thmbox}[title={Condição Suficiente de Diferenciabilidade e Cadeia de Implicações}]
  Se $f_x$ e $f_y$ existirem e forem contínuas numa vizinhança de $(a,b)$,
  então $f$ é diferenciável em $(a,b)$.

  \medskip
  \textbf{Cadeia de implicações:}
  Diferenciável $\Rightarrow$ Contínua $\Rightarrow$ Limite existe;
  mas Derivadas parciais existem $\not\Rightarrow$ Diferenciável.
\end{thmbox}

% ────────────────────────────────────────────────────────
\subsection{Exemplos Resolvidos}

\begin{exbox}{5 --- Cálculo do Gradiente}
  Determine $\nabla f$ para $f(x,y)=x^3 y + \sin(xy)$.
\end{exbox}
\begin{solbox}
  \[
    f_x = 3x^2 y + y\cos(xy), \qquad f_y = x^3 + x\cos(xy).
  \]
\end{solbox}

\begin{exbox}{6 --- Derivadas Parciais de Segunda Ordem e Hessiana}
  Para $f(x,y)=e^{x^2+y^2}$, calcule todas as derivadas parciais de
  segunda ordem.
\end{exbox}
\begin{solbox}
  $f_x = 2xe^{x^2+y^2}$, $f_y = 2ye^{x^2+y^2}$.
  \[
    f_{xx}=(2+4x^2)e^{x^2+y^2},\quad f_{yy}=(2+4y^2)e^{x^2+y^2},\quad
    f_{xy}=f_{yx}=4xye^{x^2+y^2}.
  \]
  (Note-se que $f_{xy}=f_{yx}$ pelo Teorema de Schwarz/Clairaut, uma vez
  que todas as derivadas parciais de segunda ordem são contínuas.)
\end{solbox}


% ────────────────────────────────────────────────────────
\section{Regra da Cadeia}

% ────────────────────────────────────────────────────────
\subsection{Enquadramento Teórico}

\begin{thmbox}[title={Regra da Cadeia para Funções de Várias Variáveis}]
  Sejam $z=f(x,y)$, $x=x(t)$, $y=y(t)$ diferenciáveis. Então
  \[
    \frac{dz}{dt} = \frac{\partial f}{\partial x}\frac{dx}{dt}
                  + \frac{\partial f}{\partial y}\frac{dy}{dt}.
  \]
  De modo mais geral, se $x=x(s,t)$, $y=y(s,t)$:
  \[
    \frac{\partial z}{\partial s}
      = \frac{\partial f}{\partial x}\frac{\partial x}{\partial s}
      + \frac{\partial f}{\partial y}\frac{\partial y}{\partial s}.
  \]
\end{thmbox}

% ────────────────────────────────────────────────────────
\subsection{Exemplos Resolvidos}

\begin{exbox}{7 --- Aplicação da Regra da Cadeia}
  Seja $z = x^2+y^2$, $x=\cos t$, $y=\sin t$. Calcule $\dfrac{dz}{dt}$.
\end{exbox}
\begin{solbox}
  \[
    \frac{dz}{dt} = 2x(-\sin t)+2y(\cos t) = -2\cos t\sin t + 2\sin t\cos t = 0.
  \]
  (Faz sentido: $z = \cos^2 t + \sin^2 t = 1$, constante.)
\end{solbox}
